\documentclass[10pt,letterpaper]{article}
\usepackage[utf8]{inputenc}
\usepackage[T1]{fontenc}
\usepackage[spanish,es-nodecimaldot]{babel}
\usepackage[margin=1.8cm,top=1.6cm,bottom=1.6cm]{geometry}
\usepackage{booktabs,tabularx,xcolor,hyperref,enumitem,graphicx,fancyhdr,caption}
\hypersetup{colorlinks=true,linkcolor=blue,urlcolor=blue}
\pagestyle{fancy}\fancyhead{}\fancyfoot{}\fancyfoot[C]{\thepage}\renewcommand{\headrulewidth}{0pt}
\setlength{\parskip}{.25em}\setlength{\parindent}{0pt}
\captionsetup{font=small,skip=2pt}
\setlist{nosep,leftmargin=*}
\renewcommand{\arraystretch}{1.08}

% ---------------------------------------------------------------------------
% CIFRAS PROVISIONALES: salen de results/xai/xai_summary_synthetic.json (datos
% sintéticos). Se muestran en rojo con \pv{}. Al correr src/explain.py con STEAD,
% actualizar los valores de abajo desde results/xai/xai_summary.json y redefinir
% \pv{#1} como {#1} (sin color). Cambiar también \figsuf a vacío para usar las
% figuras reales (xai_local.png, xai_global.png).
% ---------------------------------------------------------------------------
\newcommand{\pv}[1]{\textcolor{red}{#1}}
\newcommand{\figsuf}{_synthetic}

\newcommand{\vAfter}{\pv{71\,\%}}          % masa de |IG| posterior a t*
\newcommand{\vReach}{\pv{0.5 s}}           % alcance medio hacia el futuro
\newcommand{\vKeptZero}{\pv{2\,\%}}        % alertas que sobreviven sin futuro
\newcommand{\vKeptTwo}{\pv{98\,\%}}        % ... con 0.2 s de futuro
\newcommand{\vZ}{\pv{43\,\%}}\newcommand{\vN}{\pv{22\,\%}}\newcommand{\vE}{\pv{35\,\%}}
\newcommand{\vIGOcc}{\pv{0.72}}            % Spearman IG vs oclusión
\newcommand{\vIGRand}{\pv{0.95}}           % IG modelo entrenado vs aleatorio
\newcommand{\vOccRand}{\pv{0.40}}          % oclusión entrenado vs aleatorio
\newcommand{\vBase}{\pv{0.95}}             % IG línea base cero vs ruido
\newcommand{\vSteps}{\pv{0.99}}            % IG 64 vs 16 pasos
\newcommand{\vDelta}{\pv{2.6\,\%}}         % error relativo de completitud
\newcommand{\vAmp}{\pv{0.52}}              % IG vs amplitud de la señal
\newcommand{\vDelIG}{\pv{1.79}}\newcommand{\vDelOcc}{\pv{1.84}}\newcommand{\vDelRnd}{\pv{0.06}}
\newcommand{\vCostIG}{\pv{265 ms}}\newcommand{\vCostOcc}{\pv{152 ms}}
\newcommand{\vLag}{\pv{$-0.15$ s}}         % t* menos P real

\begin{document}
\begin{center}
{\Large\textbf{Laboratorio 10: Explainable AI aplicado a AlertaP}}\\[2pt]
CC3092 Deep Learning y Sistemas Inteligentes --- Universidad del Valle de Guatemala\\
Javier Eduardo España Pacheco (23361) \quad Angel Esteban Esquit Hernández (23221)\\
Noviembre de 2026
\end{center}

\textbf{Modelo explicado.} CNN 1-D residual de AlertaP (3 componentes Z/N/E, 12 s a 100 Hz) que devuelve
un logit por muestra, $f(x)[t]$, para la fase P; entrenada con STEAD y evaluada con el criterio de aviso útil del
proyecto (umbral 0.50). Se explica sin reentrenar. \pv{Las cifras y figuras en rojo son provisionales (datos
sintéticos) hasta correr \texttt{src/explain.py} con el subconjunto de prueba de STEAD.}

\section{Investigación: métodos XAI para series de tiempo}
Método $\to$ idea; requisitos; costo; limitaciones. $T$ = muestras, $n$ = pasos o muestras de perturbación.
\begin{table}[h]\centering\footnotesize
\begin{tabularx}{\textwidth}{>{\raggedright\arraybackslash}p{2.35cm}XX>{\raggedright\arraybackslash}p{2.1cm}X}
\toprule
Método & Idea & Requiere & Costo & Limitaciones\\\midrule
Oclusión por ventanas & Reemplaza una ventana por una línea base y mide la caída de la salida. & Solo la salida (caja negra); línea base y tamaño de ventana. & $T/\mathrm{paso}$ evaluaciones. & Depende de ventana y reemplazo; entradas fuera de distribución; ignora interacciones.\\
Integrated Gradients (IG) y TSR & Integra el gradiente a lo largo del camino línea base$\to x$; TSR reescala la atribución por tiempo y luego por variable. & Gradientes; línea base. & $n$ pasos adelante+atrás (TSR añade recálculos). & Sensible a la línea base; el camino atraviesa entradas irreales; puede reflejar la amplitud de la entrada.\\
TimeSHAP & Valores de Shapley (tipo KernelSHAP) sobre observaciones o eventos podados de la secuencia. & Caja negra; conjunto de fondo. & Muestreo: miles de evaluaciones. & Varianza; la poda y el fondo cambian el resultado.\\
Dynamask & Aprende una máscara de perturbación suave que minimiza el cambio de la salida. & Gradientes; optimización por instancia; regularización. & Optimización iterativa por explicación. & Hiperparámetros delicados; resultados no únicos.\\
WinIT & Importancia de cada observación sobre predicciones en horizontes futuros. & Generador de contrafactuales. & Alto (entrenar generador). & Pensado para predicción de secuencias; costoso de montar.\\
Atención (TFT) & Usa los pesos de atención como importancia. & Modelo con atención. & Casi nulo. & No aplica a una CNN; atención no equivale a causalidad.\\
\bottomrule
\end{tabularx}
\end{table}

\textbf{¿Interpretable por diseño o explicación a posteriori?} Un modelo interpretable por diseño expone su lógica
(pocos parámetros con significado); una explicación a posteriori aproxima la de un modelo opaco y puede no ser
fiel. Sí tenemos un baseline interpretable: \textbf{STA/LTA}, un cociente de energías con un umbral. En STEAD queda por
debajo de la CNN (F1 0.862 frente a 0.957; error de picking 0.59 s frente a 0.19 s).

\textbf{Propiedades formales.} IG cumple \emph{completitud}: $\sum_{t,c}\mathrm{IG}=f(x)-f(x')$ y también sensibilidad e
invariancia ante implementación; garantiza que la atribución reparte exactamente la diferencia con la línea base, no
que esa línea base sea la relevante. La oclusión no tiene garantía formal: mide el efecto de una intervención concreta. LIME
daría fidelidad local, SHAP aditividad y eficiencia.

\textbf{Parámetros de diseño que afectan la explicación.} En IG: línea base, número de pasos y muestra objetivo; en
oclusión: tamaño, paso y valor de reemplazo de la ventana; en ambos, la definición del escalar explicado.

\section{Selección del método}
Principal: \textbf{Integrated Gradients} (captum). Descartados como principal: \textbf{oclusión} (se conserva como contraste y verificación) y \textbf{TimeSHAP}.
\begin{table}[h]\centering\footnotesize
\begin{tabularx}{\textwidth}{>{\raggedright\arraybackslash}p{3.1cm}XXX}
\toprule
Criterio & IG (elegido) & Oclusión & TimeSHAP\\\midrule
Compatibilidad (CNN 1-D, 3 canales) & Total: diferenciable; atribuye por muestra \emph{y} componente. & Total, pero a nivel de ventana (0.5 s). & Posible; pensado para recurrentes y secuencias cortas.\\
Tipo de explicación & Local aditiva (completitud), agregable a global. & Local por intervención. & Local aditiva (Shapley).\\
Costo/latencia & 64 pasos: \vCostIG{} por ventana en CPU frente a $\approx$1.6 ms de inferencia. & \vCostOcc{} por ventana. & Miles de evaluaciones; el más lento.\\
Comunicación al operador & Resalta los instantes y la componente que dispararon la alerta. & Muy intuitiva (``si se borra este tramo, no hay alerta''). & Requiere explicar Shapley.\\
Limitaciones relevantes & Línea base; depende de la amplitud (\S\ref{sec:disc}). & Resolución gruesa; datos fuera de distribución. & Varianza y fondo.\\
\bottomrule
\end{tabularx}
\end{table}

\section{Implementación y parámetros}
\textbf{Escalar explicado.} Como el modelo emite un logit por muestra, se explica la \emph{decisión}: $f(x)[t^\ast]$,
donde $t^\ast$ es la primera muestra que supera el umbral tras el primer segundo de calentamiento (o el máximo del logit si
no hubo alerta). Todo se regenera con \texttt{python src/explain.py} (semilla 2026).
\begin{table}[h]\centering\footnotesize
\begin{tabularx}{\textwidth}{>{\raggedright\arraybackslash}p{3.0cm}>{\raggedright\arraybackslash}p{3.0cm}X}
\toprule
Parámetro & Valor & Justificación\\\midrule
Línea base de IG & Ceros & Las señales tienen media cero, están filtradas (1--18 Hz) y normalizadas: cero = ausencia de señal. Con una ventana de ruido como línea base la atribución temporal correlaciona \vBase{} (Spearman).\\
Pasos de IG & 64 & Error relativo medio de completitud \vDelta{}; con 16 pasos la correlación es \vSteps{}.\\
Oclusión & Ventana 0.5 s (50 muestras), paso 25, reemplazo 0 & Escala de la fase P ($\approx$ un período a 3--10 Hz más el inicio); cero coherente con la línea base de IG.\\
Casos locales & 6 de prueba & Acierto confiado, error, cercano al umbral (mínimo $|p_{\max}-0.5|$), evento de menor SNR (subgrupo), falsa alarma si existe y rechazo correcto.\\
Muestra global & 200 eventos TP (azar, semilla 2026) & Los TP muestran qué usa el modelo cuando funciona.\\
\bottomrule
\end{tabularx}
\end{table}

\begin{figure}[h]\centering
\includegraphics[width=0.82\textwidth]{../results/xai/xai_local\figsuf.png}
\caption{Explicaciones locales. Arriba de cada caso: componentes Z/N/E, probabilidad (azul) y marcas P (verde), S (naranja) y $t^\ast$ (rojo). Abajo: $|\mathrm{IG}|$ (rojo) y oclusión positiva (negro), normalizadas.}
\label{fig:local}
\end{figure}
\begin{figure}[h]\centering
\includegraphics[width=\textwidth]{../results/xai/xai_global\figsuf.png}
\caption{Explicación global sobre eventos TP: $|\mathrm{IG}|$ medio alineado a la P real, reparto por componente y prueba de borrado (fidelidad).}
\label{fig:global}
\end{figure}

\section{Interpretación en el contexto del proyecto}
\textbf{¿Señales con sentido para el dominio?} Sí: la atribución se concentra alrededor del inicio de la P (Fig.~\ref{fig:global}),
y la componente vertical Z aporta más (Z \vZ{}, E \vE{}, N \vN{}), coherente con que la P llega sobre todo polarizada en el plano
vertical-radial. \emph{Matiz:} como el campo receptivo de la red es de $\approx\pm0.9$ s, parte de esa localización es
arquitectónica y no evidencia de que el modelo ``entienda'' la P.

\textbf{Atajos y fugas buscados.} (i) Atribución sobre el ruido previo al evento (nivel de fondo como proxy de la estación): masa $\approx 0$ antes de $P-1$ s;
(ii) uso de bordes/calentamiento: el primer segundo no se alerta y no recibe atribución relevante; (iii) dependencia de la amplitud
bruta: la correlación de $|\mathrm{IG}|$ con $|x|$ es \vAmp{}, moderada; (iv) fuga entre particiones: se descartó por diseño (partición por \emph{evento}),
pero \textbf{no por estación}, así que no se descarta que el modelo aprenda rasgos de estación.

\textbf{Hallazgo central: la red no es causal.} Las convoluciones son simétricas, y \vAfter{} de la atribución cae \emph{después} de $t^\ast$ (alcance medio \vReach{}).
Si se oculta la señal posterior a $t^\ast$, solo \vKeptZero{} de las alertas se mantienen; con 0.2 s de señal futura se mantienen \vKeptTwo{}. En streaming
la alerta se retrasaría al menos $\approx$0.2 s, así que \textbf{la anticipación reportada en el informe es optimista} (la decisión hoy se toma en $t^\ast\approx P$\vLag{} usando muestras futuras).

\textbf{Errores.} \pv{(Redactar con los casos de STEAD.) En la corrida provisional el ``error'' es una alerta prematura: $t^\ast$ cae $\approx$0.7 s antes de la P y la
atribución se concentra en un pico de energía precursor; contarla como FN es correcto porque no da aviso útil. Lo habría corregido un umbral algo mayor o exigir persistencia de la probabilidad.}

\textbf{Decisiones que cambian.} (1) Evaluar y desplegar con la \emph{causalidad} explícita (ventana causal o retraso de 0.2--0.3 s) y recalcular la anticipación;
(2) añadir una regla de revisión humana cuando la atribución no se concentre en torno al inicio de la P;
(3) exigir confirmación multiestación y validar por estación antes de usar el modelo fuera de STEAD.
\textbf{Cómo se mostraría:} al operador, una forma de onda con el tramo de $\approx$1 s resaltado (oclusión, más intuitiva), la probabilidad y los segundos hasta S; el IG por componente solo al analista.
\textbf{¿Modelo interpretable?} En STEAD STA/LTA no alcanza a la CNN (arriba), así que no sustituye al modelo; se propone mantenerlo como contraste: alerta de la CNN no respaldada por STA/LTA $\Rightarrow$ revisión.

\section{Discusión y análisis}\label{sec:disc}
\textbf{Método elegido.} IG por su resolución muestral y por componente y por su completitud. Tras usarlo cambiaríamos parte de la decisión:
la correlación de IG entre el modelo entrenado y uno con pesos aleatorios es \vIGRand{}, frente a \vOccRand{} en oclusión, y con IG--oclusión de \vIGOcc{};
es decir, IG es fiel (la prueba de borrado lo confirma: caída del logit \vDelIG{} con IG, \vDelOcc{} con oclusión y \vDelRnd{} al azar con una ventana) pero poco discriminativo respecto a los pesos
aprendidos, probablemente porque su perfil temporal está dominado por la forma del campo receptivo y por la entrada (hipótesis sin verificar). Por eso reportamos \textbf{ambos} y usaríamos la oclusión como explicación principal para el operador.

\textbf{Qué reveló que las métricas no mostraban.} La dependencia del futuro: F1, precisión y error de picking no distinguen un detector causal de uno no causal, pero cambian la anticipación real.
Esto \emph{contradice} la expectativa de que la alerta se emite en el instante en que llega la P.

\textbf{Costo y viabilidad.} IG \vCostIG{} y oclusión \vCostOcc{} por ventana (CPU) frente a $\approx$1.6 ms de inferencia: unas $10^2$ veces más. Es inviable por cada ventana
en continuo, pero viable \emph{solo para alertas emitidas}, por lote o para auditoría posterior.

\textbf{Riesgos éticos abiertos.} Una explicación convincente puede dar confianza indebida a una alerta errónea; no se ha validado por estación ni en Guatemala; un falso negativo retrasa
acciones de seguridad y un falso positivo detiene operaciones. No debe usarse para accionar maquinaria, gas o alarmas públicas sin validación sismológica, redundancia y certificación,
ni en estaciones o ruido distintos a los de STEAD.

\vspace{2pt}
\textbf{Repositorio:} \url{https://github.com/AngelEsquit/PRYF-DL}\\ (sección ``Explicabilidad'' del README; \texttt{python src/explain.py}).
\end{document}
